% Options for packages loaded elsewhere
\PassOptionsToPackage{unicode}{hyperref}
\PassOptionsToPackage{hyphens}{url}
\documentclass[
  12pt,
]{ctexart}
\usepackage{xcolor}
\usepackage{amsmath,amssymb}
\setcounter{secnumdepth}{-\maxdimen} % remove section numbering
\usepackage{iftex}
\ifPDFTeX
  \usepackage[T1]{fontenc}
  \usepackage[utf8]{inputenc}
  \usepackage{textcomp} % provide euro and other symbols
\else % if luatex or xetex
  \usepackage{unicode-math} % this also loads fontspec
  \defaultfontfeatures{Scale=MatchLowercase}
  \defaultfontfeatures[\rmfamily]{Ligatures=TeX,Scale=1}
\fi
\usepackage{lmodern}
\ifPDFTeX\else
  % xetex/luatex font selection
\fi
% Use upquote if available, for straight quotes in verbatim environments
\IfFileExists{upquote.sty}{\usepackage{upquote}}{}
\IfFileExists{microtype.sty}{% use microtype if available
  \usepackage[]{microtype}
  \UseMicrotypeSet[protrusion]{basicmath} % disable protrusion for tt fonts
}{}
\makeatletter
\@ifundefined{KOMAClassName}{% if non-KOMA class
  \IfFileExists{parskip.sty}{%
    \usepackage{parskip}
  }{% else
    \setlength{\parindent}{0pt}
    \setlength{\parskip}{6pt plus 2pt minus 1pt}}
}{% if KOMA class
  \KOMAoptions{parskip=half}}
\makeatother
\usepackage{longtable,booktabs,array}
\usepackage{calc} % for calculating minipage widths
% Correct order of tables after \paragraph or \subparagraph
\usepackage{etoolbox}
\makeatletter
\patchcmd\longtable{\par}{\if@noskipsec\mbox{}\fi\par}{}{}
\makeatother
% Allow footnotes in longtable head/foot
\IfFileExists{footnotehyper.sty}{\usepackage{footnotehyper}}{\usepackage{footnote}}
\makesavenoteenv{longtable}
\setlength{\emergencystretch}{3em} % prevent overfull lines
\providecommand{\tightlist}{%
  \setlength{\itemsep}{0pt}\setlength{\parskip}{0pt}}
\usepackage{bookmark}
\IfFileExists{xurl.sty}{\usepackage{xurl}}{} % add URL line breaks if available
\urlstyle{same}
\hypersetup{
  hidelinks,
  pdfcreator={LaTeX via pandoc}}

\author{SCX}
\date{2026}

\begin{document}

\section{Proposition 1: Global Expert Ranking
Insufficiency}\label{proposition-1-global-expert-ranking-insufficiency}

\begin{quote}
全局专家排序不足定理：在异质输入空间中，不存在一个与状态无关的全局最优专家排序。
\end{quote}

\begin{center}\rule{0.5\linewidth}{0.5pt}\end{center}

\subsection{1 Statement}\label{statement}

\subsubsection{1.1 正式陈述}\label{ux6b63ux5f0fux9648ux8ff0}

设 \(\mathcal{F} = \{f_1, \dots, f_M\}\)
为一个专家集合，\(\Pi: \mathcal{X} \to \{s_1, \dots, s_K\}\)
为输入空间的一个可分划。

若存在两个状态 \(s_1, s_2 \in \mathcal{S}\) 和两个专家
\(a, b \in \{1,\dots,M\}\)，使得：

\[R_a(s_1) < R_b(s_1) \quad \text{但} \quad R_a(s_2) > R_b(s_2)\]

则不存在一个全局最优专家。

\subsubsection{1.2 定义}\label{ux5b9aux4e49}

\textbf{定义 1（全局专家排序）}：全局专家排序是一个全序关系 \(\preceq\)
在专家集合 \(\{1,\dots,M\}\) 上，满足：若 \(a \preceq b\)，则认为专家
\(a\) 整体上不劣于专家 \(b\)。

\textbf{定义 2（全局最优专家）}：专家 \(m^*\)
称为全局最优，若对任意其他专家 \(m \neq m^*\) 和任意状态
\(s \in \mathcal{S}\)，有 \(R_{m^*}(s) \leq R_m(s)\)。

\begin{center}\rule{0.5\linewidth}{0.5pt}\end{center}

\subsection{2 Formal Proof}\label{formal-proof}

\subsubsection{2.1 反证法（一）}\label{ux53cdux8bc1ux6cd5ux4e00}

假设存在全局最优专家 \(m^*\)。

由全局最优的定义，对所有 \(s \in \mathcal{S}\)：

\[R_{m^*}(s) \leq R_a(s) \quad \text{且} \quad R_{m^*}(s) \leq R_b(s)\]

但在状态 \(s_1\) 处：\(R_{m^*}(s_1) \leq R_a(s_1) < R_b(s_1)\)，故
\(m^*\) 在 \(s_1\) 上优于 \(b\)。

在状态 \(s_2\) 处：\(R_{m^*}(s_2) \leq R_b(s_2) < R_a(s_2)\)，故 \(m^*\)
在 \(s_2\) 上也优于 \(a\)。

交叉条件 \(R_a(s_1) < R_b(s_1)\) 且 \(R_a(s_2) > R_b(s_2)\) 意味着：

\[R_a(s_1) - R_b(s_1) < 0 < R_a(s_2) - R_b(s_2)\]

由于 \(R_{m^*}(s_1) \leq R_a(s_1)\) 且
\(R_{m^*}(s_2) \leq R_b(s_2) < R_a(s_2)\)，我们有：

\[R_{m^*}(s_1) \leq R_a(s_1) < R_b(s_1)\]
\[R_{m^*}(s_2) \leq R_b(s_2) < R_a(s_2)\]

但这不产生矛盾。矛盾需要额外的条件。因此使用更精确的构造法。

\subsubsection{2.2 构造法（反例）}\label{ux6784ux9020ux6cd5ux53cdux4f8b}

设 \(\mathcal{X} = [0, 1]\)，输入分布 \(P_X \sim U[0, 1]\)。使用平方损失
\(\ell(\hat{y}, y) = (\hat{y} - y)^2\)。

定义状态 \(s_1 = [0, 0.5)\), \(s_2 = [0.5, 1]\)，概率权重
\(P_X(s_1) = 0.9, P_X(s_2) = 0.1\)。

构造函数：

\begin{longtable}[]{@{}llll@{}}
\toprule\noalign{}
专家 & \(R(s_1)\) & \(R(s_2)\) & 全局 \(R\) \\
\midrule\noalign{}
\endhead
\bottomrule\noalign{}
\endlastfoot
\(f_1\) & 1 & 100 & \(0.9 \times 1 + 0.1 \times 100 = 10.9\) \\
\(f_2\) & 2 & 5 & \(0.9 \times 2 + 0.1 \times 5 = 2.3\) \\
\end{longtable}

全局上 \(R(f_1) = 10.9 > 2.3 = R(f_2)\)，专家 \(f_2\) 全局更好。但在状态
\(s_1\) 中 \(R_1(s_1) = 1 < 2 = R_2(s_1)\)，即全局更优的专家 \(f_2\)
在状态 \(s_1\) 中反而更差。

因此，不存在能在两个状态上同时最优的单个专家。\(\square\)

\subsubsection{2.3
修正证明（状态条件最优集）}\label{ux4feeux6b63ux8bc1ux660eux72b6ux6001ux6761ux4ef6ux6700ux4f18ux96c6}

设 \(\mathcal{M} = \{1,\dots,M\}\)。定义状态条件最优集：

\[\mathcal{M}^*(s) = \arg\min_{m \in \mathcal{M}} R_m(s)\]

若 \(a \in \mathcal{M}^*(s_1)\) 且 \(b \in \mathcal{M}^*(s_2)\)，且
\(a \neq b\)，则
\(|\bigcap_{s \in \mathcal{S}} \mathcal{M}^*(s)| \leq 1\)。若
\(|\mathcal{S}| \geq 2\)
且存在至少两个不同的最优专家，则不存在单个专家在所有状态上都是最优的。

由条件 \(R_a(s_1) < R_b(s_1)\) 得 \(a \in \mathcal{M}^*(s_1)\) 或
\(b \notin \mathcal{M}^*(s_1)\)；由 \(R_a(s_2) > R_b(s_2)\) 得
\(b \in \mathcal{M}^*(s_2)\) 或 \(a \notin \mathcal{M}^*(s_2)\)。

若同时有 \(a \in \mathcal{M}^*(s_1)\) 和
\(b \in \mathcal{M}^*(s_2)\)，则当 \(a \neq b\) 时得证。若 \(a = b\)，则
\(R_a(s_1) < R_a(s_1)\)（矛盾），故 \(a \neq b\) 必然成立。

因此，\(\mathcal{M}^*(s_1) \neq \mathcal{M}^*(s_2)\)，即全局最优专家不存在。\(\square\)

\begin{center}\rule{0.5\linewidth}{0.5pt}\end{center}

\subsection{3 Corollary: Co-monotonicity
Condition}\label{corollary-co-monotonicity-condition}

全局专家排序 \(R_m \leq R_{m'}\) 蕴含状态级排序
\(R_m(s) \leq R_{m'}(s)\) 对所有 \(s\) 成立，当且仅当对任意两个专家
\(f_m, f_{m'}\)，差值函数

\[\Delta_{m,m'}(x) = \ell(f_m(x), f^*(x)) - \ell(f_{m'}(x), f^*(x))\]

与输入分布 \(P_X\) 的似然比 \(dP_X/dP_X(\cdot|s)\)
处处共单调。即存在函数 \(\psi_{m,m'}\) 使得：

\[\Delta_{m,m'}(x) \cdot \frac{dP_X}{dP_X(\cdot|s)}(x) \geq 0, \quad \forall x \in \mathcal{X}, \forall s \in \Pi\]

当 \(\mathcal{S}\) 连续时（如 \(\mathcal{S} = [0,1]\)），命题退化为：若
\(\exists s_1, s_2\) 使得 \(R_a(s_i)\) 与 \(R_b(s_i)\) 交叉，则不存在
Lebesgue-几乎必然的全局最优专家：

\[\mu\left(\left\{s \in \mathcal{S} : m^*(s) = \arg\min_m R_m(s)\right\}\right) < 1\]

其中 \(\mu\) 是 \(\mathcal{S}\) 上的某参考测度（如输入分布 \(P_X\) 的
push-forward）。

\begin{center}\rule{0.5\linewidth}{0.5pt}\end{center}

\subsection{4 Implications for SCX}\label{implications-for-scx}

\begin{enumerate}
\def\labelenumi{\arabic{enumi}.}
\item
  \textbf{必要性论证}：这一命题直接论证了 SCX
  框架的\textbf{必要性}而非充分性。它说明：

  \begin{itemize}
  \tightlist
  \item
    任何全局聚合指标（平均精度、平均 AUC、F1 宏观平均）都会丢失信息
  \item
    在交叉发生的数据区域，必须引入状态条件建模
  \item
    交叉的频次决定了状态粒度的需求：交叉越多，需要的状态越多
  \end{itemize}
\item
  \textbf{专家选择必须状态条件化}：全局 accuracy
  不足以描述专家能力。专家路由应基于状态条件风险 \(R_m(s)\) 而非全局风险
  \(R_m\)。
\item
  \textbf{与 Dawid-Skene 模型的关系}：Dawid-Skene
  假设标注者可靠性是全局常数（\(\pi^{(j)}_{kl}\) 不依赖 \(x\)）。SCX
  允许可靠性随状态 \(s\) 变化。当存在状态 \(s_1, s_2\) 使得 Dawid-Skene
  全局混淆矩阵无法同时预测 \(s_1\) 和 \(s_2\)
  上的表现时，状态条件建模是必要的。这正是命题一的本质。
\end{enumerate}

\begin{center}\rule{0.5\linewidth}{0.5pt}\end{center}

\subsection{参考文献}\label{ux53c2ux8003ux6587ux732e}

\begin{enumerate}
\def\labelenumi{\arabic{enumi}.}
\tightlist
\item
  SCX 核心框架数学分析. \texttt{01\_SCX\_核心框架\_数学分析.md} Section
  2.1
\item
  SCX 数学基础与证明建构. \texttt{05\_数学根源与证明.md} Proposition 1
\item
  Arrow, K. J. (1951). \emph{Social Choice and Individual Values}.
\item
  Dawid, A. P. \& Skene, A. M. (1979). Maximum likelihood estimation of
  observer error-rates. \emph{Applied Statistics}.
\end{enumerate}

\end{document}
